% !TEX root = ../../../main.tex
\section{Inverse function theorem}
\label{sec:analysis-i:inverse-function-theorem}
% Source: images/LEC06/IMG_9351--IMG_9352, inverse derivative calculation.

\begin{pendingreview}
The inverse differentiation rule used for roots and logarithms is a
one-dimensional statement. Continuity of a strictly monotone inverse was
established in \autoref{thm:analysis-i:continuous-inverse}.
\begin{theorem}[Derivative of an inverse]
  \label{thm:analysis-i:inverse-derivative}
  Let $f:I\to J$ be a continuous strictly monotone bijection between
  open intervals, and let $g=f^{-1}$. If $f$ is differentiable at
  $x\in I$ and $f'(x)\ne0$, then $g$ is differentiable at $y=f(x)$ and
  \[
    g'(y)=\frac1{f'(x)}=\frac1{f'(g(y))}.
  \]
\end{theorem}
\begin{proof}
  Restrict $f$ to a closed subinterval of $I$ containing $x$ in its
  interior. By \autoref{thm:analysis-i:continuous-inverse} and its
  decreasing analogue, its inverse, and hence $g$, is continuous near $y$.
  For $k\ne0$ sufficiently small, set $h=g(y+k)-g(y)$. Injectivity
  gives $h\ne0$, and continuity of $g$ gives $h\to0$ as $k\to0$.
  Also $k=f(x+h)-f(x)$, so
  \[
    \frac{g(y+k)-g(y)}{k}
      =\frac1{\bigl(f(x+h)-f(x)\bigr)/h}
      \longrightarrow\frac1{f'(x)}.
  \]
\end{proof}
This difference-quotient argument establishes that the inverse derivative
exists before applying the chain rule.
\end{pendingreview}
Differentiating $g(f(x))=x$ now gives the identity
\[
  g'(f(x))f'(x)=1.
\]
% The pending argument fills a local justification, not a general
% higher-dimensional inverse function theorem.
